\documentclass[11pt]{article}
\usepackage{amsmath,amssymb}
\usepackage[margin=1in]{geometry}
\setlength{\parindent}{1.2em}
\title{Rewritten Chain of Thought for the Solution to the\\ Two-Minus-Sector Water-Wave Amplitude $A_n$}
\author{}
\date{}
\begin{document}
\maketitle
\vspace{-3em}
\begin{abstract}
\noindent This document contains a rewritten summary of the chain of thought for the
closed-form solution of the tree-level $n$-point on-shell amplitude $A_n$ for one-dimensional
deep-water surface gravity waves in the two-minus sector
$\sigma=(-1,-1,+1,\dots,+1)$, in PDF form.
\end{abstract}

\section*{Rewritten Chain of Thought}

The object is the tree amplitude $A_n$ produced by the Berends--Giele recursion in
\texttt{OnShellBG.m}, evaluated on the resonant manifold
\[
\sum_{i=1}^{n}\omega_i = 0,
\qquad
\sum_{i=1}^{n}\sigma_i\,\omega_i^2 = 0,
\]
in the two-minus sector $\sigma=(-1,-1,+1,\dots,+1)$. The dispersion law $\omega_i^2=g\,|k_i|$
fixes the momenta as $k_i=\sigma_i\omega_i^2/g$. I set $g=1$ until the very end, where
dimensional analysis restores it. Legs $1,2$ are the two negative-$\sigma$ legs and legs
$3,\dots,n$ the positive ones. The hint promises a piecewise homogeneous polynomial in the
frequencies, and the task is to produce it.

Before generating any data it pays to ask where non-polynomiality could possibly enter. The
recursion is assembled from rational kernels, a propagator $-i/(\omega_S^2/|k_S|-g)$, and the
magnitude $\mathrm{mag}[k]=|k|$. The only non-rational operation anywhere is that absolute
value, and it is applied exclusively to \emph{momenta}, never to frequencies. An external
momentum is $k_i=\sigma_i\omega_i^2$, so $|k_i|=\omega_i^2$ is unambiguous. The recursion,
however, forms \emph{internal} momenta
\[
k_S=\sum_{i\in S}\sigma_i\,\omega_i^2
\]
on subsets $S$, and $|k_S|=\big|\sum_{i\in S}\sigma_i\omega_i^2\big|$ is genuinely
sign-dependent. This is the only place a piecewise structure can come from: the chambers must
be the regions on which every $k_S$ keeps a fixed sign. I keep this in reserve and start
measuring.

A handful of exact rational evaluations settle the gross features. Every $A_n$ comes out purely
imaginary, and rescaling all frequencies by $\lambda$ multiplies the value by
$\lambda^{2(n-2)}$ --- at $n=5$ the ratio is $64=2^{6}$, at $n=6$ it is $256=2^{8}$. So
\[
A_n=-\,i\,P_n,\qquad P_n\ \text{real, homogeneous of degree } 2(n-2),
\]
and that degree is the first hard constraint on any guess.

The $n=4$ evaluation refuses to return a number; it is $0\cdot\infty$. The kinematic solver,
which fixes two frequencies from the conservation laws, has no freedom at $n=4$ but to land on
$\omega_4=-\omega_2$ (and $\omega_1=-\omega_3$), forcing the internal momentum
$k_{\{2,4\}}=-\omega_2^2+\omega_4^2=0$. A vanishing internal momentum sends the propagator to
zero while a kernel develops a $0/0$, and the symbolic engine cannot resolve the product.
Choosing a different pair of solved legs only moves the zero to another channel. The conclusion
is structural rather than accidental: the entire $n=4$ two-minus variety lies on a wall where
some $k_S=0$. I set $n=4$ aside as a boundary case to be recovered later, and work at
$n\ge 5$, where the interior is non-empty.

Permuting the two minus legs, or permuting the plus legs among themselves, preserves both
conservation laws and so maps on-shell points to on-shell points. Direct evaluation shows $A_n$
is invariant under all of these: a generic point and all $2\cdot(n-2)!$ of its images return one
and the same value. Thus $A_n$ is symmetric under $S_2\times S_{n-2}$.

This symmetry seems to hand me the answer. On the constraint surface the two minus-leg
symmetric functions reduce to those of the plus legs: writing $P_k$ for the $k$-th elementary
symmetric polynomial of $\omega_3,\dots,\omega_n$, conservation gives $\omega_1+\omega_2=-P_1$
and $\omega_1\omega_2=P_2$. Hence any $S_2\times S_{n-2}$-symmetric polynomial of degree
$2(n-2)$, restricted to the surface, is a weighted-homogeneous polynomial in $P_1,\dots,P_{n-2}$
alone --- for $n=5$ that is exactly seven monomials. I generate points, set up the $7\times7$
linear system, and solve.

It does not fit. Not approximately: with exact rational points and a full-rank system the
residual on held-out points is of order $10^{11}$. Suspecting a degenerate sample, I widen and
diversify the points; the residual stays enormous. This is a genuine contradiction. The function
is symmetric (checked directly), homogeneous of a definite degree (checked directly), and the
hint says polynomial --- yet it is provably not the restriction of any symmetric polynomial,
since such a restriction would have to lie in the seven-dimensional span that the data refute.

There is only one way out, and it is the chamber idea I had parked. A symmetric \emph{function}
need not be a single symmetric polynomial if it is merely \emph{piecewise} polynomial: the
global symmetry can be realised by gluing together pieces that are individually \emph{not}
symmetric, with the symmetry group permuting the pieces. The walls are the sign changes of the
internal momenta $k_S$. Within one chamber every absolute value resolves to a fixed sign, so the
amplitude is a single rational function and --- by the hint --- a polynomial; but across a wall
$|k_S|$ flips to $-k_S$ and the polynomial can change. The piece in any one chamber carries no
obligation to be symmetric, only to glue continuously to its neighbours. The fit failed because
it demanded a global symmetry that no single piece possesses.

To see one piece cleanly I collapse to a line: fix $\omega_3,\omega_4$ at rational values, let
$\omega_2=t$ vary, and let the solver give $\omega_1,\omega_5$ as functions of $t$. A symbolic
attempt drowns --- the nested absolute values of rational functions generate expressions large
enough to exhaust memory --- so instead I evaluate the exact rational $A_5(t)$ at many rational
$t$ inside a single chamber and reconstruct the function. Multiplying by the apparent denominator
$\omega_2+\omega_3+\omega_4=13+t$ clears it to a polynomial,
\[
A_5(t)\,(13+t)=208\,t^6+2128\,t^5=16\,t^5\,(13t+133),
\]
and since the solver returns $\omega_1=-(13t+133)/(13+t)$, the denominator cancels against the
numerator factor, leaving the clean monomial
\[
A_5=-16\,\omega_1\,\omega_2^5
\]
on this chamber. The factor of $\omega_1$ was hidden inside the apparent pole; the result depends
on the plus legs not at all. Checking against scattered points of the same chamber --- wildly
different plus-leg data, the identical monomial --- confirms it. And the chamber is recognisable:
it is the region in which $\omega_2$ is the unique smallest frequency, equivalently the one in
which $k_{\{2\}}$ is the \emph{only} negative subset momentum.

The monomial $-16\,\omega_1\omega_2^5$ is wildly asymmetric in the two minus legs, exactly as it
should be: the chamber ``$\omega_2$ smallest'' is not invariant under
$\omega_1\leftrightarrow\omega_2$ --- that swap leaves the chamber --- so the piece living on it
is free to be asymmetric. The same monomial recurs at every $n$ in the smallest chamber, with a
clean power of two for a coefficient; evaluating at $n=6,7$ gives
\[
A_n=-\,2^{\,n-1}\,\omega_1\,\omega_2^{\,2n-5}\qquad(\omega_2\ \text{smallest}),
\]
whose exponents $1$ and $2n-5$ sum to the required $2(n-2)$.

Now I let $\omega_2$ grow and watch the amplitude cross walls. As $\omega_2^2$ rises past each
partial sum $q_S=\sum_{i\in S}\omega_i^2$ of plus-leg squares, the subset momentum
$k_{\{2\}\cup S}=-\omega_2^2+q_S$ changes sign. Reconstructing the polynomial on each successive
chamber and differencing neighbours, the jump across the wall $\omega_2^2=q_S$ is always
\[
\Delta_S=(-1)^{|S|+1}\,2^{\,n-1}\,\omega_1\omega_2\,(\omega_2^2-q_S)^{\,n-3},
\]
which vanishes on the wall (the function is continuous) and switches on only beyond it. Starting
from the smallest chamber --- where only $S=\varnothing$ is active and the piece is
$-2^{n-1}\omega_1\omega_2(\omega_2^2)^{n-3}=-2^{n-1}\omega_1\omega_2^{2n-5}$ --- and adding the
jump for every wall already crossed, the amplitude assembles into a single closed form,
\[
A_n=-\,i\,2^{\,n-1}\,\omega_1\omega_2
\sum_{\substack{S\subseteq\{3,\dots,n\}\\ \omega_2^2>q_S}}
(-1)^{|S|+1}\,(\omega_2^2-q_S)^{\,n-3},
\]
the sum running over precisely those subsets whose wall has been crossed. Equivalently, with the
truncated power $(x)_+=\max(x,0)$, it is the spline
\[
A_n=-\,i\,2^{\,n-1}\,\omega_1\omega_2
\sum_{S\subseteq\{3,\dots,n\}}(-1)^{|S|+1}\big[(\omega_2^2-q_S)_+\big]^{\,n-3}.
\]
A reassuring instance: in the chamber $\omega_3^2+\omega_4^2<\omega_2^2$ at $n=5$, where all of
$\varnothing,\{3\},\{4\},\{3,4\}$ are active, the four truncated powers combine, the $\omega_2^4$
and $\omega_2^2$ terms cancel, and only $-32\,\omega_1\omega_2\,\omega_3^2\omega_4^2$ survives ---
again a clean monomial, now in the four small legs.

Two facts make the formula consistent rather than merely fitted. First, the exponent is $n-3$
while the number of plus legs is $n-2$, and the $(n-2)$-fold alternating sum of a polynomial of
degree $n-3$ vanishes identically,
\[
\sum_{S\subseteq\{3,\dots,n\}}(-1)^{|S|}\,(x-q_S)^{\,n-3}=0\qquad\text{for all }x,
\]
which I verified symbolically for $n=4,\dots,7$. This is why the spline closes up: the jumps
across a full nested set of walls sum to zero, leaving nothing at the far end. Second, the same
identity disposes of an apparent ambiguity. I wrote the threshold with $\omega_2^2$, but
$S_2$-symmetry demands it not matter whether I threshold on $\omega_2^2$ or $\omega_1^2$. On shell
$\omega_1^2+\omega_2^2=q_{\{3,\dots,n\}}$ and $q_{S^c}=q_{\{3,\dots,n\}}-q_S$; fed through the
vanishing identity these relate the truncated sum built on $\omega_1^2$ to the one built on
$\omega_2^2$, and numerically the two agree at every point and in every chamber, including those
where $\omega_1$ is the smaller leg. The formula is genuinely well defined on the symmetric
variety, not an artefact of labelling.

The boundary case now falls out. At $n=4$ the variety is forced onto the wall
$\omega_2^2=\omega_4^2$, so the Berends--Giele evaluation sits exactly on a discontinuity of the
absolute value and returns $0\cdot\infty$. But the spline is continuous there, with the definite
value $-i\,8\,\omega_1\omega_2^3$, and approaching the wall through nearby off-resonant points the
recursion converges to precisely this number. The $0\cdot\infty$ is a removable singularity whose
value the closed form supplies.

Finally the coupling. Restoring $g$, the amplitude is not invariant: comparing $g=1$ with $g=2$
shows $A_n\propto g^{-(n-3)}$. This is forced once one sees that the natural variable inside the
bracket is the magnitude $|k_j|=\omega_j^2/g$ rather than $\omega_j^2$; written in terms of
$|k_j|$ the factors of $g$ cancel, and only the explicit $g^{-(n-3)}$ from re-homogenising
remains,
\[
\boxed{\;
A_n=-\,i\,\frac{2^{\,n-1}}{g^{\,n-3}}\,\omega_1\omega_2
\sum_{S\subseteq\{3,\dots,n\}}(-1)^{|S|+1}\big[(\omega_2^2-q_S)_+\big]^{\,n-3},
\qquad q_S=\sum_{i\in S}\omega_i^2.
\;}
\]

What remains is to be sure the guess is not merely a good fit on the line where it was found.
Calling the recursion on one representative of every chamber reachable by sweeping $\omega_2$ ---
at $n=7$ these range from the one-term smallest chamber up to chambers whose polynomial carries
sixteen truncated-power terms --- the closed form reproduces the exact rational Berends--Giele
value with identically zero residual; the same holds at hundreds of random points and at explicit
configurations with $\omega_1$ smallest, for $n=5,6,7$, together with the limiting agreement at
$n=4$. Zero residual in exact arithmetic is stronger than any tolerance: the piecewise polynomial
\emph{is} the amplitude.

In hindsight the structure is natural. The combination
$\sum_S(-1)^{|S|}\big[(x-q_S)_+\big]^{m}$ is the standard inclusion--exclusion / truncated-power
building block of a one-dimensional spline with knots $\{q_S\}$; here the evaluation point is
$x=\omega_2^2$, the knots are the partial sums of plus-leg squares, and the two minus legs supply
the prefactor $\omega_1\omega_2$. The amplitude is, up to that prefactor and the power of two, a
spline in the single variable $\omega_2^2$ against the knot set $\{q_S\}$; its polynomial pieces
are the kinematic chambers, and the vanishing of the full alternating sum --- degree one below
the number of knots --- is exactly the condition that the pieces glue into one globally
symmetric, globally continuous function.

\end{document}
